\scanpage{111}
\mspage{55}

\margin{Sous l'hyp. ``suite exacte d'homotopie'', mais avec fibres pas néc. \unsure{conn.}}

on aura une factorisation can. de $X\longrightarrow Y$ en
\[
(14)\qquad X\longrightarrow Y'\longrightarrow Y
\]
avec $Y'$ étale fini ou pro-étale fini sur $Y$ et $\mathcal E_X^{\unsure{\prime}}\longrightarrow\mathcal E_{Y'}$ est maintenant épi, \deleted{\unsure{donc}} \unsure{cette|avec} suite exacte univ. d'homotopie bien sûr. On aura des iso can
\[
(15)\qquad \mathcal E_{X/Y}^T\xrightarrow{\sim}\mathcal E_{X/Y'}^T
\]
qui peut donc se calculer comme
\[
(16)\qquad \mathcal E_{X\times_{Y'}T}^T\simeq\mathcal E_{X/Y}^T
\]
qu'on peut noter $\mathcal E_{X_T}^T$, mais en faisant attention que maintenant $X_T$ désigne non plus $X\times_YT$ (qui va être \unsure{disconnexe} \inserted{si $Y'\longrightarrow Y$ pas iso}) mais $X\times_{Y'}T$.

Bien sûr, à iso. can. près on peut \deleted{\unsure{en}} remplacer, pour le calcul de $\mathcal E_{X_T}^T$, $T$ par $\widetilde Y$ pour ce qui est du $X_T$, et par $\widetilde X$ pour ce qui est de \unsure{l'ancrage|l'origine} dans $X_{\widetilde Y}$,

\margin{110}
